% ICASSP 2027 submission -- DRAFT
% Deadline: 16 September 2026.  4 pages content + 1 page references.
% Every number in this file is traceable to a run record under results/runs/.
% Placeholders are marked TODO.

\documentclass[10pt,conference]{IEEEtran}
\IEEEoverridecommandlockouts
\usepackage{amsmath,amssymb,graphicx,booktabs,url,balance}
\usepackage[hidelinks]{hyperref}

\begin{document}

\title{Is the Band the Biomarker?\\
A Validity Protocol for Band-Specific EEG Classification Claims}

% TODO: confirm whether ICASSP 2027 is anonymous. Historically it is NOT.
\author{\IEEEauthorblockN{TODO Author}
\IEEEauthorblockA{TODO Affiliation\\
TODO email}}

\maketitle

\begin{abstract}
Resting-state EEG studies routinely report that one frequency band is the
strongest discriminator for a clinical label. We observe that such claims are
seldom checked for well-posedness before modelling begins, and we propose
\emph{BandCheck}: five inexpensive tests that ask whether a band-specific claim
can be true at all. The tests examine passband support, the accuracy
attributable to the cross-validation split alone, whether a five-parameter
signal-quality baseline matches the model, whether a narrow sub-band reproduces
the whole result, and whether excising a suspected contaminant changes anything.
We apply the protocol to a published result reporting 99.60\% accuracy for major
depressive disorder using the 30--100\,Hz band. Four of five checks fail. The
nominal band extends 20\,Hz beyond the recordings' 80\,Hz acquisition low-pass;
the split rule alone accounts for $0.106\pm0.013$ accuracy over three seeds; a five-scalar baseline with no
neural content matches an eleven-layer CNN ($0.873$ vs $0.847\pm0.009$, McNemar
$p\geq0.625$ across seeds); and a single uniform third of the band reproduces the full-band
predictions ($r=0.99$, 98\% agreement). We reproduce $0.847\pm0.009$
subject-level accuracy, not $0.9960$. We release the protocol as a tool; the checks
add minutes to a linear baseline and a small multiple of one model fit to a deep
one.
\end{abstract}

\begin{IEEEkeywords}
EEG, biomarkers, validity, reproducibility, spectral analysis, depression
\end{IEEEkeywords}

% ======================================================================
\section{Introduction}

A recurring result form in clinical EEG reads: \emph{model $M$ applied to band
$B$ classifies condition $C$ with accuracy $a$, and $B$ is the most effective
band}. The form is attractive because it appears to localise a biomarker in
frequency. It is also fragile in ways that are cheap to test and rarely tested.

Three failure modes recur. First, the nominal band may not be supported by the
data: acquisition filters bound the spectrum, and a band specified beyond that
bound contains no signal. Second, the reported accuracy may depend on the
cross-validation split rather than on the signal; epoch-level splits allow a
model to recognise individuals rather than the condition. Third, the
discriminative information may be non-neural --- muscle, electrode impedance,
mains --- in which case the band is incidental.

None of these requires a replication study to detect. Each is a short
computation that can be run before any modelling. We package them as
\textbf{BandCheck}, a five-check validity protocol, and demonstrate it on a
recent high-profile result.

\textbf{Contributions.} (i) A protocol of five checks that establish whether a
band-specific claim is well-posed, with a released reference implementation.
(ii) An application to a published 99.60\% MDD classification result, in which
four of five checks fail. (iii) A record of six alternative explanations that we
tested and rejected, which we argue should accompany any such analysis.

We emphasise that the protocol tests \emph{well-posedness}, not truth. A claim
passing all five checks may still be wrong; a claim failing one is not yet
meaningful enough to be either.

% ======================================================================
\section{The BandCheck protocol}

Let a claim specify a band $B=[f_\mathrm{lo},f_\mathrm{hi}]$, a model $M$, a
dataset $D$ of recordings from $N$ subjects, and a reported accuracy $a$.

\subsection{V1: Passband support}
Estimate the usable spectral edge $f_\mathrm{edge}$ from (a) the acquisition
filter settings recorded in the file header and (b) the frequency above which
the median power spectral density falls to the noise floor. The dead fraction is
$\phi = \max(0, f_\mathrm{hi}-f_\mathrm{edge})/(f_\mathrm{hi}-f_\mathrm{lo})$.
We flag $\phi>0.02$ and fail $\phi\geq0.15$. Header and empirical estimates are
reported separately; disagreement is itself informative.

\subsection{V2: Protocol delta}
Fit $M$ twice, changing only the split rule: epochs assigned at random, versus
subjects assigned wholly to train or test. The difference $\Delta_p$ is accuracy
attributable to the protocol. We flag $\Delta_p\geq0.03$ and fail
$\Delta_p\geq0.10$.

\subsection{V3: Artifact control}
Fit a logistic regression on five per-epoch scalars computed within $B$:
relative 60--100\,Hz power, relative 45--55\,Hz power, kurtosis, zero-crossing
rate, and RMS amplitude. None encodes a neurophysiological construct; together
they describe signal quality. Compare to $M$ under identical subject-independent
evaluation using McNemar's test on paired subject-level decisions. If the
baseline matches or exceeds $M$, there is no evidence that $M$ exploits neural
information.

\subsection{V4: Sub-band localisation}
Partition $B$ into $k$ uniform sub-bands (we use $k=3$) and fit $M$ on each.
Report the accuracy of the best sub-band, its correlation with the full-band
predicted probabilities, and their decision agreement. A narrow sub-band that
reproduces the full-band predictions means the claim is mis-specified in width.
A uniform partition is used deliberately: it is not tuned to the answer.

\subsection{V5: Excision control}
Band-stop a suspected contaminant (by default the mains region) and refit. If
accuracy is unchanged, that contaminant is not the mechanism. This check most
often returns a null, and that is its value: it prevents a plausible but wrong
mechanism from being asserted.

% ======================================================================
\section{Case study}

\subsection{Target claim and data}
Anik \emph{et al.}~\cite{anik2024} report 99.60\% accuracy for MDD versus healthy
control using an eleven-layer 1-D CNN (Ex-1DCNN) on the gamma band, defined as
30--100\,Hz, with 15\,s epochs and tenfold cross-validation, and conclude that
gamma is the most effective band. The data are the public HUSM
recordings~\cite{mumtaz2017}: 19 scalp electrodes, 256\,Hz, eyes-closed and
eyes-open resting conditions.

We reimplemented Ex-1DCNN from the published layer table. Three data properties
required decisions, all of which we report because they change the denominator:
the release contains a duplicated recording (two copies of one eyes-open file);
one MDD subject has only task data and therefore no resting recording; and seven
subjects lack one of the two resting conditions. Our analysis uses the 63
subjects with resting data (33 MDD, 30 HC; 119 recordings), against the 34/30
stated in~\cite{anik2024}.

We also note a reproducibility detail: the regularisation parameter
$\lambda=0.02$ stated in~\cite{anik2024} collapses our reimplementation to a
constant output when applied as decoupled weight decay, at every learning rate
tested. We use $10^{-4}$. We report this neutrally; the published $\lambda$ is
evidently a different formulation.

\subsection{Report}
Table~\ref{tab:bandcheck} gives the BandCheck report for the claim. All results
use 15\,s epochs, tenfold subject-wise cross-validation, and a training budget of
40 epochs; the budget is stated because it moves the numbers (Sec.~\ref{sec:sens}).

\begin{table}[t]
\caption{BandCheck report for the 30--100\,Hz MDD claim of~\cite{anik2024}.
$n=63$ subjects, Ex-1DCNN.}
\label{tab:bandcheck}
\centering
\small
\begin{tabular}{@{}llp{4.6cm}@{}}
\toprule
\textbf{Check} & \textbf{Status} & \textbf{Result} \\
\midrule
V1 Passband  & \textbf{FAIL} & 29\% of the band (20\,Hz) lies above the 80\,Hz usable edge \\
V2 Protocol  & \textbf{FAIL} & split rule alone accounts for $0.106\pm0.013$ (2/3 seeds exceed threshold) \\
V3 Artifact  & \textbf{FAIL} & 5-scalar baseline matches the CNN: $0.873$ vs $0.847\pm0.009$, McNemar $p\geq0.625$ \\
V4 Sub-band  & \textbf{FAIL} & 30--53.3\,Hz alone matches or exceeds the full band in 4/4 runs; $r=0.99$, 98\% agreement \\
V5 Excision  & PASS$^\dagger$ & removing 45--55\,Hz changes accuracy by $-0.005\pm0.033$; sign flips across seeds \\
\midrule
\multicolumn{3}{c}{\textbf{Verdict: not well-posed} (4/4 runs)} \\
\multicolumn{3}{@{}p{7.5cm}@{}}{\footnotesize $^\dagger$V5's categorical label is seed-dependent (PASS/PASS/FLAG/FLAG); the quantity is zero within noise.} \\
\bottomrule
\end{tabular}
\end{table}

\subsection{V1: a fifth of the band was never recorded}
The EDF headers report \texttt{HP:0.5Hz LP:80Hz} for all 120 resting files, with
no difference between groups. The nominal band therefore extends 20\,Hz beyond
the acquisition low-pass. This is confirmed empirically: 70--100\,Hz yields a
median MDD/HC log-power ratio of $-0.01$ and mean per-bin AUC $0.537$, and
100--127\,Hz yields $0.550$ --- chance, as expected of empty spectrum.

\subsection{V3: five scalars suffice}
The artifact baseline reaches $0.873$ subject-level accuracy against
$0.847\pm0.009$ for Ex-1DCNN (McNemar $p\geq0.625$ in every run; 95\% CIs
$[0.778,0.952]$ and $[0.746,0.921]$ at seed~0). We state this as a match, not a win: with
$N=63$ the smallest detectable difference is approximately $0.095$
(Sec.~\ref{sec:power}), and the observed gap is $0.016$.

The same five scalars reach $0.81$--$0.88$ in every band we tested --- delta,
beta, gamma, the mains region, and gamma with the mains region removed --- a
spread of $0.070$ against a majority-class baseline of $0.524$. Band choice is
therefore close to irrelevant for this discriminant, which is difficult to
reconcile with a claim that one band is uniquely effective.

\subsection{V4: the effect is 15--23\,Hz wide, not 70}
The best uniform third, 30--53.3\,Hz, is 30--53.3\,Hz in all four runs and
matches or exceeds the full band in each ($0.873$ vs $0.857$ at seed~0), with
predicted probabilities correlating at $r=0.99$ and 98\% decision agreement. A hand-specified 30--45\,Hz band --- reported as a secondary,
post-hoc analysis, and measured at the reduced 25-epoch budget --- gives
identical subject-level accuracy to the full band ($0.8413$ both, McNemar
$p=1.00$, $r=0.979$, 93.7\% agreement). Everything above
45\,Hz, three-quarters of the nominal band, is inert.

\subsection{What the discriminative signal is}
Fitting an aperiodic $1/f^{\chi}$ model over 2--45\,Hz (mains excluded), the
exponent differs strongly between groups (MDD $2.15$, HC $1.32$; AUC $0.776$,
$p=1.8\times10^{-4}$), as does the offset (AUC $0.859$,
$p=1.1\times10^{-6}$). These two numbers alone give $0.794\pm0.084$
subject-level cross-validated accuracy.

Removing the fitted aperiodic component does not remove the group difference
(all-band CV $0.858\rightarrow0.841$), so a periodic component also contributes.
We therefore do not claim to have identified the mechanism. What the protocol
establishes is narrower and, we think, more useful: the reported effect is not a
30--100\,Hz phenomenon, is not driven by mains, and is matched by a
signal-quality baseline.

% ======================================================================
\section{Alternative explanations tested and rejected}

Table~\ref{tab:rejected} lists hypotheses we advanced and then rejected. We
report them because a negative claim is only as credible as the alternatives
ruled out, and because each rejection is itself a measurement.

\begin{table}[t]
\caption{Hypotheses tested and rejected.}
\label{tab:rejected}
\centering
\small
\begin{tabular}{@{}lp{4.3cm}@{}}
\toprule
\textbf{Hypothesis} & \textbf{Rejected by} \\
\midrule
Subject leakage explains the result & protocol delta only $0.106$; subject-ID probe $13.4\times$ chance \\
Temporalis (jaw) EMG & topography is posterior, not temporal; $>$55\,Hz at chance \\
Controls have more mains contamination & raw 50\,Hz peak ratio AUC $0.568$, $p=0.36$ \\
Groups recorded with different filters & EDF headers identical \\
The mains region drives gamma & excising 45--55\,Hz costs $-0.016$ \\
Aperiodic slope explains everything & flattening changes CV by $-0.017$ \\
\bottomrule
\end{tabular}
\end{table}

% ======================================================================
\section{Sensitivity and statistical power}

\subsection{Seed and budget sensitivity}\label{sec:sens}
We ran the full protocol four times: three seeds at a 40-epoch training budget
and one at 25 epochs. \textbf{The overall verdict is \emph{not well-posed} in
all four runs}, and the V1, V3 and V4 outcomes are identical in all four.

Two checks have seed-dependent \emph{labels}. The V2 quantity is
$0.106\pm0.013$, straddling our $0.10$ threshold: two of three seeds fail, one
flags. The V5 quantity is $-0.005\pm0.033$ and its sign flips across seeds,
swinging the label between pass and flag; in subject counts the effect is one to
two people out of 63.

We regard this as the strongest available argument for a recommendation we would
otherwise have had to assert: \textbf{report the quantity, not the categorical
verdict}. A threshold that sits inside the seed-to-seed spread carries no
information. Note also that the V5 result is \emph{strengthened} by this
analysis --- an effect centred on zero with flipping sign is stronger evidence
that the mains region is not the mechanism than a single small negative value
would be.

Budget variation is smaller than seed variation: the 25-epoch V2 delta
($0.114$) lies inside the 40-epoch seed range $[0.091, 0.116]$. No verdict
depends on the training budget.

\subsection{Power}\label{sec:power}
With $N=63$ subjects, McNemar cannot reach $p<0.05$ until at least six subjects
net change classification, corresponding to an accuracy difference of $0.095$.
No subject-level model comparison below this threshold is detectable in this
dataset. We therefore avoid model-vs-model claims, report confidence intervals
throughout, and rest the analysis on epoch-level contrasts, where $n$ is in the
thousands, and on structural findings that require no significance test.

\subsection{Threshold calibration}
The thresholds in Sec.~II ($\phi\geq0.15$, $\Delta_p\geq0.10$, $\pm0.02$ for V5,
a half-width criterion for V4) are reasoned defaults, not calibrated ones. They
are exposed as parameters. We recommend reporting the underlying quantity rather
than the categorical verdict. Sec.~\ref{sec:sens} shows this is not a
hypothetical concern: two of our five checks change label across seeds while
every underlying quantity stays stable.

% ======================================================================
\section{Discussion}

The protocol is deliberately cheap. V1 requires reading a file header. V3 and V5
are minutes of computation. The most expensive check, V4, is $k$ additional model
fits. Against the cost of a replication study, or of a biomarker claim that does
not hold, this is negligible.

We make no claim about the biology of depression, and we do not claim
that~\cite{anik2024} is methodologically unusual. The evaluation convention it
follows is widespread, which is precisely why a mechanical check is worth
having. Our own analysis produced six hypotheses that the data rejected; we
regard that ratio as typical rather than exceptional, and as an argument for
running the checks early.

\textbf{Limitations.} Single dataset; $N=63$; the mechanism behind the
30--45\,Hz effect is unresolved, with residual EMG, the aperiodic difference,
and a genuine low-gamma effect all remaining consistent with our observations.
Distinguishing them requires an independent cohort. The thresholds are
uncalibrated. Our reimplementation may differ from the original in unstated
details; we release code and per-run records so that any discrepancy can be
located.

% ======================================================================
\section{Conclusion}

Band-specific EEG classification claims can be checked for well-posedness before
they are believed. Applied to a published 99.60\% result, four of five checks
fail: a fifth of the nominal band was never recorded, the split rule accounts
for nine accuracy points, a five-parameter signal-quality baseline matches an
eleven-layer network, and a single uniform third of the band reproduces the
full-band predictions. We reproduce $0.847\pm0.009$ subject-level accuracy. The protocol
and all run records are released.

\section*{Reproducibility}
Code, configurations and a timestamped record of every run reported here ---
including the command, git commit and package versions --- are available at
\url{https://github.com/TODO/neuro-or-noise}. Table~\ref{tab:bandcheck} is
produced by a single command.

\balance
\begin{thebibliography}{9}
\bibitem{anik2024}
I.~A. Anik, A.~H.~M. Kamal, M.~A. Kabir, S.~Uddin, and M.~A. Moni,
``A robust deep-learning model to detect major depressive disorder utilizing EEG
signals,'' \emph{IEEE Trans. Artif. Intell.}, vol.~5, no.~10, pp.~4938--4947,
Oct. 2024.

\bibitem{mumtaz2017}
W.~Mumtaz, L.~Xia, S.~S.~A. Ali, M.~A.~M. Yasin, M.~Hussain, and A.~S. Malik,
``Electroencephalogram (EEG)-based computer-aided technique to diagnose major
depressive disorder,'' \emph{Biomed. Signal Process. Control}, vol.~31,
pp.~108--115, 2017.

\bibitem{whitham2007}
E.~M. Whitham \emph{et al.}, ``Scalp electrical recording during paralysis:
quantitative evidence that EEG frequencies above 20\,Hz are contaminated by
EMG,'' \emph{Clin. Neurophysiol.}, vol.~118, no.~8, pp.~1877--1888, 2007.

\bibitem{muthu2013}
S.~D. Muthukumaraswamy, ``High-frequency brain activity and muscle artifacts in
MEG/EEG: a review and recommendations,'' \emph{Front. Hum. Neurosci.}, vol.~7,
p.~138, 2013.

\bibitem{geirhos2020}
R.~Geirhos \emph{et al.}, ``Shortcut learning in deep neural networks,''
\emph{Nature Mach. Intell.}, vol.~2, pp.~665--673, 2020.

\bibitem{donoghue2020}
T.~Donoghue \emph{et al.}, ``Parameterizing neural power spectra into periodic
and aperiodic components,'' \emph{Nature Neurosci.}, vol.~23, pp.~1655--1665,
2020.

% TODO: add 3-5 further EEG-MDD papers that use epoch-level evaluation, to
% support the claim that the convention is widespread rather than specific
% to [1]. This citation set is load-bearing for the paper's framing.

\end{thebibliography}

\end{document}
